一维质量 \(m\) 的量子粒子在双吸引点势 \(V(x)=-g[\delta(x-a)+\delta(x+a)]\) 中运动，\(g>0,a>0\)，令 \(\alpha=mg/\hbar^2\)。考察能量 \(E=-\hbar^2\kappa^2/(2m)<0\) 的束缚态，波函数在两点连续，但导数满足分布势导致的跳跃条件；忽略自旋。
(1) 推导两个\(\delta \)点处的波函数导数跳跃条件与奇偶宇称分解。
(2) 分别建立偶宇称和奇宇称束缚态的超越方程，并证明偶态始终存在。
(3) 求奇束缚态出现的临界条件，检查两点重合 \(a\to0\) 的极限。
(4) 在 \(\alpha a\gg1\) 的远隔势阱极限给出两束缚态的指数小能量劈裂。
